\documentclass[12pt]{article}
\usepackage[utf8]{inputenc}
\usepackage[T1]{fontenc}
\usepackage{hyperref}
\title{Varhadi-Nagpuri Language Resources for K-12 Curriculum Development}
\author{Compiled by Claude Sonnet 5 (claude-sonnet-5) for Andr\'as Kornai}
\date{}
\begin{document}
\maketitle
\section{Introduction}
Varhadi (Varhadi-Nagpuri, ISO 639-3 vah, Glottocode varh1239) is spoken mainly in Vidarbha, Maharashtra. Wikipedia describes it as a dialect of Marathi and lists about 7 million speakers; it is described as Vulnerable in the UNESCO Atlas. My searches found no Varhadi-specific textbook, primer, dictionary or grammar. This dossier is therefore short.
\section{Primary Educational Resources}
No Varhadi-medium primers or textbooks were found in search results. Instruction in the region is presumably in standard Marathi or Hindi.
\section{Linguistic and Reference Materials}
\begin{itemize}
\item Varhadi dialect, Wikipedia (overview, grammar notes such as the case ending \emph{le}): \url{https://en.wikipedia.org/wiki/Varhadi_dialect}
\item Marathi-English dictionary entry for the term Varihadi (a dictionary entry for the word, not a Varhadi lexicon): \url{https://www.wisdomlib.org/definition/varihadi}
\item Wiktionary entry for the word Varhadi: \url{https://en.wiktionary.com/wiki/Varhadi}
\end{itemize}
\section{Religious and Community Resources}
None found.
\section{Digital and Audio Resources}
None found.
\section{Curriculum Development Context}
Varhadi is treated as a Marathi dialect, so a curriculum would need original documentation work, for example by local universities or the Central Institute of Indian Languages. These are leads only; I did not verify any such materials.
\end{document}